\lecture{1}{Sun 30 Jan 2022 15:42}{abstract error Secretarhy}
\section{Abstract Error Secretary Problem with Examination}
\subsection{Problem Formulation}
$n$ items arrive in random order. Let $i$ denote the item that arrives in the round $i$ .  
The items are assigned absolute ranks $r_i$ . $r$ is totally ordered wrt i.  
Define $d_i$ to be the player's decision at round $i$ . It equals one of the following: $\text{Accept}, \text{Decline}, \text{Examine}$  
\begin{itemize}
	\item $d_i = \text{Accept} \implies \text{Terminates}$  
	\item $d_i = \text{Accept} \land \forall j\in [n], r_j < r_i \iff \text{Wins}$  
\end{itemize}
Define $O_i$ to be the items that the player has \textbf{observed} by round i, and $E_i$ to be the items that player has \textbf{examined} by round i.  
\begin{itemize}
	\item $O_1 = [1], E_1 = \emptyset$  
	\item $E_i \subseteq O_i$  
	\item $d_i = \text{Decline} \implies O_{i+1} = O_i \cup \{i+1\}, E_{i+1} = E_i$  
	\item $d_i = \text{Examine} \implies O_{i+1} = O_i , E_{i+1} = E_{i} \cup \{i\}$  
\end{itemize}

\begin{definition}[Indistinguishable Items]
$i \approx j$ means that $i$ and $j$ are in-distinguishable given that neither of them have been examinated.
	\begin{itemize}
 \item (reflexivity) $i \approx i$
  \item(symmetry) $i \approx j \iff j \approx i$
  \item(downward closure) $i \approx j , r_i < r_j \implies \forall k, l [r_k, r_l \in [r_i, r_j] \implies r_k \approx r_l ]$
\end{itemize}
\end{definition}   

Define $N_i := \{j: j \approx i\}$ to be the \textit{neighbour} of $i$. Note that $1\leq |N_i| \leq n$ .  
Define $\Theta_i$ to be the relations known to the player by round i. 
\begin{itemize}
	\item $\Theta_k$ defines a partial order $\le_{k}$ on $O_k$ , such that $i \le_{k} j \iff (i,j) \in \Theta_k$  
	\item $\Theta_1 = \{(1,1)\}$  
	\item $i \le_k j \implies i \le_{k+1} j$  
	\item $\forall i,j \in O_k[ i \not\approx j \implies i \le_k j \lor j \le_k i]$  
	\item ($E_k$ totally ordered) $\forall i,j \in E_k[i \le_k j \lor j \le_k i]$ .  
\end{itemize}
Summary: For an incoming item $i$ , upon first observation, the player can compare $r_i$ and $r_j$ for $\forall j \in O_i \setminus N_i$ . Upon the second observation, the player can compare the examined item with \textbf{previously examined} items (which can be in $N_i$ ).  


\subsection{Algorithmic Representations}


\begin{algorithm}[H]
\caption{Player's algorithm for the problem}
\label{algo1} 
\begin{algorithmic}[1]
\For{$i \in [n]$}
	\State{$\alpha \gets$ the current item to be decided}
	\If{$i < r$} \Comment Phase 1
		\If{$\exists \beta\in O_i \ldotp \alpha <_i \beta$}\Comment $\alpha$ is \emph{not} a maximum
			\State{Decline this item}
		\ElsIf{$\alpha \in E_i$}
			\State{Decline this item}
		\Else 
			\State{Examine this item}
		\EndIf
	\ElsIf{$i < n$}\Comment Phase 2
		\If{$\exists \beta\in O_i \ldotp \alpha <_i \beta$}
			\State{Decline this item}
		\Else				\If{$\forall \beta\in O_i \ldotp \beta \leq_i \alpha$} \Comment $\alpha$ is the absolute maximum
				\State{Accept this item}
			\ElsIf{$\alpha \in E_i$}
				\State{Decline this item}
			\Else
				\State{Examine this item}
			\EndIf
		\EndIf
	\Else
		\State{Accept this item}
	\EndIf

\EndFor

\end{algorithmic}
\end{algorithm}




\begin{algorithm}[H]
\caption{Game Process}
\label{algo1} 
\begin{algorithmic}[1]


\State{$[n], r_i$ and $\approx$ relation are given according to definition}


\State{$O_1 = \{1\}$}
\State{$E_1 = \emptyset$}
\State{$\alpha_1 = 1$}
\State{$\Theta_1 = \emptyset$}



\For{$i \in [n]$}
	\State{$d_{i} \gets $ player's decision}
	\If{$d_i = \mathrm{Accept}$}
		\State Game terminates with reward $v(\alpha_i)$
			
	\ElsIf{$d_i = \mathrm{Decline}$}
		\State{$O_{i+1} \gets O_i\cup \{i+1\}$}
		\State{$E_{i+1} \gets E_i$}
		\State{$\alpha_{i+1} \gets i+1$}
		
	\ElsIf{$d_i = \mathrm{Examine}$}
		\State{$O_{i+1} \gets O_i$}
		\State{$E_{i+1} \gets E_i\cup \{i\}$}
		\State{$\alpha_{i+1} \gets \alpha$}
	\EndIf
	\State{$\Theta_{i+1} \gets 
		\{(\beta, \gamma) \in O_{i+1}^2: r_\beta \leq r_\gamma, \beta \not\approx \gamma\} \cup
		\{(\beta, \gamma) \in E_{i+1}^2: r_\beta \leq r_\gamma\}$}
	
	
\EndFor

\end{algorithmic}
\end{algorithm}


\section{Relaxation of the abstract error SP problem}
In the previous definition, $i$ and $j$ need to both have been examinated, for the player to compare them precisely.
Now we relax the problem, needing only one examination between $i$ and $j$ to obain precise comparison.

We give the expected value of $\text{ALG}$ on the relaxed model. Let $\alpha$ denote the selected item.

\begin{equation}
  \operatorname{E}\left[\text{OPT}\right]=m
\end{equation}

\begin{equation}
  \operatorname{E}[ \text{ALG} ]=\sum_{t=1}^m t\operatorname{Pr}\left[{v(\alpha)=t}\right]
\end{equation}

\begin{align}
  \operatorname{Pr}\left[{v(\alpha)=t}\right] &= \sum_{i=r}^n
  \operatorname{Pr}\left[{\alpha = i | v(i)=t}\right]
  \operatorname{Pr}\left[{v(i)=t}\right]\\&+
  \operatorname{Pr}\left[{\text{reach }n \cap \text{not accept }n}\right] \operatorname{Pr}\left[{v(n)=t}\right]\\&+
  \operatorname{Pr}\left[{\text{reach }n-1 \cap \text{ED }n-1}\right]\operatorname{Pr}\left[{v(n-1)=t}\right]
\end{align}

The two added terms are cases where the items are exhausted. If $n$ is reached then it must be selected. If $n-1$ is examined then it must be selected also.

\begin{equation}
  \operatorname{Pr}\left[{v(i)=t}\right]=1/m, \forall i\forall t
\end{equation}

\begin{align}
  \operatorname{Pr}\left[{\alpha = i | v(i)=t}\right] &= \operatorname{Pr}\left[{\text{reach }i|\max_{j<r}v(j)\leq t}\right]\operatorname{Pr}\left[{\max_{j<r}v(j)\leq t}\right]\\&=
  \sum_{s=1}^t \operatorname{Pr}\left[{\text{reach }i|\max_{j<r}v(j)= s}\right]\operatorname{Pr}\left[{\max_{j<r}v(j)= s}\right]
\end{align}

\begin{equation}
	\label{eqn:exhaust1}
	\operatorname{Pr}\left[{\text{reach }n \cap \text{not accept }n}\right]
	=\sum_{s=t+1}^m
		\operatorname{Pr}\left[{\text{reach }n|\max_{j<r}v(j)= s}\right]\operatorname{Pr}\left[{\max_{j<r}v(j)= s}\right]
\end{equation}

\begin{equation}
	\label{eqn:exhaust2}
	\operatorname{Pr}\left[{\text{reach }n-1 \cap \text{ED }n-1}\right]
	=\sum_{s=t+1}^m
		\textbf{1}_{s\leq t + \delta}
		\operatorname{Pr}\left[{\text{reach }n-1|\max_{j<r}v(j)= s}\right]\operatorname{Pr}\left[{\max_{j<r}v(j)= s}\right]
\end{equation}

\begin{align}
\operatorname{Pr}\left[{\max_{j<r}v(j)= s}\right] 
	&= \operatorname{Pr}\left[{\max_{j<r}v(j)\leq s}\right]-\operatorname{Pr}\left[{\max_{j<r}v(j)\leq s-1}\right]\\
	&= (\frac{s}{m})^{r-1}-(\frac{s-1}{m})^{r-1}
\end{align}

\begin{align}
  &\operatorname{Pr}\left[{\text{reach }i|\max_{j<r}v(j)= s}\right] \\
  &\qquad = \sum_{(a_{1\dots k}) \in\{1,2\}^k,\sum a_j = i-r} \,\prod_{j=1}^k 
  \textbf{1}_{\delta<t}
  \left(\textbf{1}_{a_j = 1}\frac{s-\delta-1}{m}
  +\textbf{1}_{a_j = 2}\frac{\delta}{m}
  \right)
  +\textbf{1}_{\delta\geq t}\frac{t-1}{m}
  \\
  \begin{split}
        &\qquad = \sum_{x=0}^{\lfloor\frac{i-r}{2}\rfloor} 
  \left(
  		\textbf{1}_{\delta<t}\left(\frac{\delta}{m}\right)^x\left(\frac{s-\delta-1}{m}\right)^{i-r-2x}
  		+\textbf{1}_{\delta \geq t} \left(\frac{t-1}{m}\right)^x 0^{i-r-2x}
  \right)\\
    &\qquad\qquad \times\frac{(i-r-x)!}{x!(i-r-2x)!}
  \end{split}
\end{align}

Explanation: There are $i-r$ steps to go from $r$ to $i$. Each decline will go 1 step, and each examine+decline will go 2 steps. We calculate the probability to take each possible path and then sum them up. In the second equation, $x$ is the number of items examined going from $r$ to $i$.
